\documentclass[10pt]{article}
\usepackage[margin=0.78in]{geometry}
\usepackage{amsmath,amssymb,amsthm}
\usepackage{graphicx}
\usepackage{microtype}
\usepackage{xcolor}
\usepackage{hyperref}
\hypersetup{colorlinks=true,linkcolor=blue!55!black,urlcolor=blue!55!black,citecolor=blue!55!black}
\newtheorem{theorem}{Theorem}
\newtheorem{corollary}[theorem]{Corollary}
\newcommand{\N}{\mathbb N}
\newcommand{\B}{\mathcal B}
\newcommand{\NT}{\mathcal N_T}
\title{Lacunary gaps answer Shkarin's positive-density-subset question}
\author{Candidate full-solution packet for arXiv:1710.07901}
\date{11 August 2026}

\begin{document}
\maketitle

\begin{center}
\fbox{\parbox{0.93\linewidth}{\textbf{Status:} candidate full positive answer, likely valid.  The frequently hypercyclic vector already constructed in Grivaux's proof has a half-space return set contained in dyadic blocks with infinitely many factor-four gaps.  Every subset having a natural density must consequently have density zero.  In particular, $T=2B$ on $\ell^2(\N_0)$ gives an explicit operator answering Question~3.3.}}
\end{center}

\section{Source question}

Sophie Grivaux, \emph{Frequently hypercyclic operators with irregularly
visiting orbits}, arXiv:1710.07901; J. Math. Anal. Appl. 462 (2018),
542--553~\cite{Grivaux}.  On page 11 of the source PDF, Question~3.3 asks:

\begin{quote}
Does there exist an operator $T\in\B(X)$ admitting a frequently hypercyclic
vector $x\in X$ such that, for some non-empty open subset $U\subseteq X$,
$\NT(x,U)$ contains no subset $A$ of $\N$ admitting a positive density?
\end{quote}

Here
\[
 \NT(x,U)=\{n\geq1:T^nx\in U\},\qquad
 \operatorname{dens}A=\lim_{N\to\infty}\frac{\#(A\cap[1,N])}{N}
\]
when the latter limit exists.  Thus ``positive density'' in the question
means positive natural asymptotic density, not merely positive upper density.

\begin{figure}[h]
 \centering
 \includegraphics[width=\linewidth]{figures/open_problem_crop.png}
 \caption{The complete source statement, page 11, Question 3.3.}
\end{figure}

\section{Proof Intuition}

Grivaux distributes the centers used to build her frequently hypercyclic
vector only in dyadic blocks $[2^j,2^{j+1})$ whose indices satisfy
$j\equiv0$ or $2\pmod 5$.  The functional used to detect returns can see only
a short interval around each center, and the construction places every
center far enough inside its dyadic subblock that this entire interval stays
in the same dyadic block.  Hence the detected return set is completely empty
on the two consecutive blocks with indices $5q+3$ and $5q+4$.

For a subset $A$ of this return set, its count is therefore unchanged while
the observation scale grows from approximately $2^{5q+3}$ to
$2^{5q+5}$, a factor of four.  If $A$ had a density $\delta$, the two endpoint
subsequences would force both $\delta$ and $\delta/4$ as the same limit.
Thus $\delta=0$.  This elementary consequence of the source construction is
stronger than the ``return set has no density'' conclusion recorded there.

\section{The strengthened theorem}

We use exactly the hypotheses of Theorem~1.2 in the published source.  A
vector $x_0$ is \emph{bicyclic} for $T$ if it admits a bilateral orbit
$(T^nx_0)_{n\in\mathbb Z}$ whose linear span is dense.

\begin{theorem}\label{thm:main}
Let $T\in\B(X)$ satisfy the following conditions from Grivaux's
Theorem~1.2:
\begin{enumerate}
 \item $T$ has a bicyclic vector $x_0$;
 \item the associated bilateral-orbit series is unconditionally convergent;
 \item there are $0\ne x_0^*\in X^*$ and a finite set
 $F\subseteq[-d,d]\cap\mathbb Z$ such that
 $\langle x_0^*,T^nx_0\rangle=0$ for $n\notin F$.
\end{enumerate}
Then $T$ admits a frequently hypercyclic vector $x$ such that, for the
non-empty open half-space
\[
 U_0=\{y\in X:\operatorname{Re}\langle x_0^*,y\rangle>0\},
\]
the return set $\NT(x,U_0)$ contains no subset of positive natural density.
In fact, the same $x$ has this property for every non-empty open
$V\subseteq U_0$, and every such $\NT(x,V)$ has no natural density.
\end{theorem}

\begin{proof}
We audit the particular vector constructed in the proof of
Theorem~1.2 of~\cite{Grivaux}.  To keep the return set distinct from the
source's center sets, write those center sets as $C_s$ here.  The construction
starts with dyadic blocks
\[
 I_j=[2^j,2^{j+1})\cap\N
\]
and pairwise disjoint subintervals $I_j^{(s)}\subseteq I_j$.  Choose $p\geq1$
so that
\begin{equation}\label{eq:margin}
 2^{s+1+p}\geq2^{s+1}+2d+1\qquad(s\geq1).
\end{equation}
Inside $I_j^{(s)}$, the set $L_j^{(s)}$ consists of selected integer centers
$k$ satisfying
\begin{equation}\label{eq:center}
 \operatorname{dist}\bigl(k,\N\setminus I_j^{(s)}\bigr)
 \geq2^{s+1+p}.
\end{equation}
The source then defines
\begin{equation}\label{eq:Cs}
 C_s=\coprod_{\substack{j\geq2s+p+2\\j\in J}}L_j^{(s)},
 \qquad J=5\N\cup(5\N+2).
\end{equation}

The vector built from these $C_s$ is frequently hypercyclic by the source's
Lemma~2.1.  Its Lemma~2.2 gives the following exact decomposition of the
half-space return set $R:=\NT(x,U_0)$:
\begin{equation}\label{eq:Rdecomp}
 R=\coprod_{s\geq1}\coprod_{k\in C_s}I_{k,s}^+,
 \qquad
 I_{k,s}^+\subseteq[k-(2^s+d),\,k+(2^s+d)]\cap\N.
\end{equation}
If $k\in L_j^{(s)}$, then \eqref{eq:margin}--\eqref{eq:center} imply
\[
 2^s+d<2^{s+1+p}
 \quad\Longrightarrow\quad
 I_{k,s}^+\subseteq I_j^{(s)}\subseteq I_j.
\]
Together with \eqref{eq:Cs} and \eqref{eq:Rdecomp}, this proves the key
containment
\begin{equation}\label{eq:keycontain}
 R\subseteq\bigcup_{j\in J}I_j.
\end{equation}

For every sufficiently large integer $q$, neither $5q+3$ nor $5q+4$ belongs
to $J$.  Consequently \eqref{eq:keycontain} makes $R$ empty on
\[
 [2^{5q+3},2^{5q+5})\cap\N.
\]
Let $A\subseteq R$, and put
\[
 a(N)=\#(A\cap[1,N]),\qquad
 N_q=2^{5q+3}-1,
 \qquad M_q=2^{5q+5}-1.
\]
The gap gives the exact equality $a(M_q)=a(N_q)$.  If $A$ has a natural
density $\delta$, then
\[
 \delta
 =\lim_{q\to\infty}\frac{a(M_q)}{M_q}
 =\lim_{q\to\infty}
   \frac{a(N_q)}{N_q}\frac{N_q}{M_q}
 =\frac{\delta}{4},
\]
because $N_q/M_q\to1/4$.  Hence $\delta=0$.  No subset of $R$ can have
positive density, proving the first assertion.

Finally let $\varnothing\ne V\subseteq U_0$ be open.  Then
$\NT(x,V)\subseteq R$, so it too contains no positive-density subset.
Since $x$ is frequently hypercyclic,
$\underline{\operatorname{dens}}\NT(x,V)>0$.  Were its natural density to
exist, it would therefore be positive, a contradiction.  Thus
$\NT(x,V)$ has no natural density as claimed.
\end{proof}

\section{An explicit operator}

\begin{corollary}\label{cor:shift}
Question~3.3 has an affirmative answer with $X=\ell^2(\N_0)$ and
$T=2B$, where $B$ is the unilateral backward shift.
\end{corollary}

\begin{proof}
Let $x_0=e_0$ and define its bilateral orbit by
\[
 x_{-m}=2^{-m}e_m\quad(m\geq1),
 \qquad x_n=0\quad(n\geq1).
\]
Then $Tx_n=x_{n+1}$ for every $n\in\mathbb Z$, and the negative-index
vectors span all coordinate vectors.  The bilateral-orbit series is
unconditionally convergent (indeed, absolutely in norm), because
$\sum_{m\geq0}\|2^{-m}e_m\|<\infty$.  With
$x_0^*(y)=y_0$, one has
$\langle x_0^*,x_n\rangle=0$ for every $n\ne0$.  Thus all hypotheses of
Theorem~\ref{thm:main} hold with $F=\{0\}$.
\end{proof}

The same verification works for $\lambda B$ with $|\lambda|>1$ on
$\ell^p(\N_0)$, $1\leq p<\infty$, or on $c_0$, after replacing
$2^{-m}$ by $\lambda^{-m}$.  The single explicit Hilbert-space example in
Corollary~\ref{cor:shift} already answers the source question in full.

\section{Verification, limitations, and novelty status}

The proof uses no numerical approximation.  The accompanying exact checker
confirms the integer endpoint convention and the convergence of the endpoint
ratio to $1/4$ on sample scales; it is a sanity check, not a substitute for
the symbolic proof.  The decisive source audit is pages 7--8: Lemma~2.2 gives
\eqref{eq:Rdecomp}, and page 8 gives the margin condition and
$J=5\N\cup(5\N+2)$.  The explicit $2B$ orbit was also checked coordinate by
coordinate.  The separate verification report records these checks.

A bounded novelty search through 11 August 2026 covered the run registry and
attempt indexes and web/arXiv searches for the exact question wording,
arXiv identifier, title, ``positive density subset'', ``irregularly visiting
orbit'', and close variants.  It found the source and its journal version,
but no later paper explicitly resolving Question~3.3.  The proof should
therefore be treated as a candidate new observation, with moderate novelty
confidence pending a specialist citation search.

The theorem does not answer whether \emph{every} frequently hypercyclic
operator has such a vector, nor whether one vector has irregular returns to
every bounded open set or every open ball.  It does give the natural class of
all non-empty open subsets contained in one detecting half-space.  Human
review is recommended, especially of the one-line support containment
\eqref{eq:keycontain} and the interpretation of ``density'' as the natural
density defined in the source.

\begin{thebibliography}{9}
\bibitem{Grivaux}
S. Grivaux,
\emph{Frequently hypercyclic operators with irregularly visiting orbits},
J. Math. Anal. Appl. 462 (2018), 542--553,
\href{https://doi.org/10.1016/j.jmaa.2018.02.020}{doi:10.1016/j.jmaa.2018.02.020};
\href{https://arxiv.org/abs/1710.07901}{arXiv:1710.07901}.
\end{thebibliography}

\end{document}
